% A283 -- Thermodynamics in FFGFT: a comprehensive overview
% Johann Pascher -- ORCID 0009-0000-6518-4064
% Chapter source for A-series; Wrapper: wr_standalone_A4/A283_Thermodynamics_FFGFT_En.tex
% Status: August 2026

\providecommand{\qedsq}{\raisebox{0.03em}{\framebox[0.62em]{\rule{0pt}{0.42em}}}}

% ===================== ABSTRACT =====================
\begingroup\small\noindent
\textbf{Abstract.}
Thermodynamics is anchored in FFGFT at five levels --- not as a
postulate but as a consequence of geometry.
This document collects the distributed individual results, arranges
them into a unified chain of argument, and identifies where the levels
interact and where they remain strictly separate.
A closing section concedes that both frameworks --- expanding space and
fractal inward refinement --- are models that interpret the same
observations in structurally opposite ways.
\par\endgroup

\bigskip

\begingroup\small\noindent
\textbf{Layer markers.}
\textbf{[STIPULATION]} declared starting point ---
\textbf{[B]} algebraically/mathematically derived from declared foundations ---
\textbf{[K]} numerically computed from corpus constants ---
\textbf{[Q]} external primary source ---
\textbf{[S]} plausibility sketch ---
\textbf{[H]} open research question.
\par\endgroup

\bigskip

% ===================== 1 =====================
\section{Why a comprehensive document}

The thermodynamic content of FFGFT emerged across several documents:
Doc.~157 treats the arrow of time, A271--A282 the Landauer accounting,
Doc.~313 the cosmological thermodynamics.
Each document is self-contained; none provides an overview of how the
parts relate to one another.

\textbf{[B]} Three questions remain unanswered without a comprehensive
document:
(1) Where do the levels touch, and where do they stay separate?
(2) Which statements depend on the $\xi$ scheme, which are independent?
(3) What is the epistemic status of each statement?
This document answers these questions and introduces no new claims.

% ===================== 2 =====================
\section{Level 1 --- The arrow of time is geometric (Doc.~157)}

\subsection{The standard problem}

All fundamental equations of physics are time-reversal invariant.
The second law ($\Delta S \ge 0$) is statistical in nature and shifts
the problem onto the unexplained low initial entropy
(Penrose: $\sim 10^{-10^{123}}$).

\subsection{The FFGFT answer}

\textbf{[B]} Three geometric arguments interlock (Doc.~157):

\textbf{Argument 1: Algebraic asymmetry.}
$\tilde{T} \cdot m = 1$ forces $\tilde{T} > 0$ because mass is
positive definite. Time reversal would require $m \to -m$ ---
physically impossible.

\textbf{Argument 2: Topological winding direction.}
The vacuum field winds with $\dot\theta = m = 1/\tilde{T}$ along
the $T^4$. This chirality is topologically protected.

\textbf{Argument 3: Fractal dimensional asymmetry.}
$D_f = 3 - \xi_0 \approx 2{.}9999$ breaks the exact mirror symmetry;
path integrals for forward and backward propagation differ by
$O(\xi_0)$.

\subsection{Consequence}

\textbf{[B]} The logic is inverted: $\Delta S \ge 0$ in FFGFT is a
\emph{consequence} of the geometry, not a postulate. The
\enquote{Past Hypothesis} becomes unnecessary --- low initial entropy
is geometrically forced, not an initial condition.

% ===================== 3 =====================
\section{Level 2 --- Landauer as phase-space accounting (A271)}

\subsection{The phase-space formulation}

\textbf{[STIPULATION]} A physical system occupies a region in phase
space of size $W$ (measured in units of $h^f$). The region exists
independently of how one names or interprets it.

\textbf{[B]} A map that reduces several source regions $\{A, B\}$ onto
a target region $Z$ shrinks the volume:
$\Delta S_\text{system} = -k_B \ln(W_{A+B}/W_Z) \le 0$.
Liouville's theorem forbids volume shrinkage in a closed system.
A heat bath must absorb the volume:
\[
  \boxed{Q \;\ge\; k_B\,T\,\ln\!\left(\frac{W_\text{before}}{W_\text{after}}\right)}
\]
with $\ln 2$ as the binary special case.

\subsection{Three separations}

\textbf{[B]}
(1) \emph{Bit $\ne$ physical region.} Shannon and Boltzmann entropy
coincide only when the description maps exactly onto the physical
microstates with equal probability.
(2) \emph{Content neutrality.} Thermodynamics counts regions and does
not read them. Bit \enquote{true} and bit \enquote{false} cost
identically.
(3) \emph{Three levels strictly separate.} Construction / state
transition / description. Landauer costs arise at levels~1 and~2,
not at~3.

% ===================== 4 =====================
\section{Level 3 --- Phase-space accounting universally (A280--A282)}

\textbf{[B]} A280: The region comes first. A glass poured into the
ocean --- the volume does not vanish, it distributes over
$1{.}4 \times 10^9$\,km³. The price: the bath must absorb the volume,
heat flows. The counting base is distinguishability (orthogonality),
not energy; absolute base discretisation: $h^f$.

\textbf{[B]} A281: Landauer's own restriction --- logically irreversible
is a description property, not a physical one. Sufficient for heat is
solely the non-bijectivity of the region map. Bennett's result: every
logically irreversible function can be executed as a sequence of
bijective region maps; costs attach to \emph{forgetting}, not to the
operation.

\textbf{[B]} A282: The token computer as a limiting case. The accounting
holds exactly; the Landauer bound does \emph{not} apply there
(engineering costs dominate). The transition to the Brownian computing
sphere connects continuously with Landauer's domain of applicability.

\textbf{[H]} Open question: minimal coherence cost for quantum coherence
independent of error correction quality (A271, §\,7.4).

% ===================== 5 =====================
\section{Level 4 --- Cosmological thermodynamics (Docs.~312/313)}

\subsection{Thermal history as geometric condition}

\textbf{[K]} Doc.~313: Three phases on the observable half-cycle
$(\tau_c/2 \approx 46\,\text{Gyr})$:

\begin{center}
\small
\begin{tabular}{llll}
\toprule
Phase & Range & $\theta/\pi$ & Character \\
\midrule
Radiation domination & $z > 3000$ & $> 0{.}90$ &
  $m \to m_\text{min}$, slowest clocks \\
Transition & $z \approx 3000$--$1090$ & $0{.}90$--$0{.}979$ &
  recombination, scatter wall \\
Matter domination & $z < 3000$ & $< 0{.}90$ &
  $m(\theta)$ grows, structure formation \\
\bottomrule
\end{tabular}
\end{center}

\subsection{The five-step chain}

\textbf{[K]} Complete derivation of $\Omega_m^*$ from $\xi_0$ alone:
\begin{center}
\small
\begin{tabular}{clc}
\toprule
& Step & Result \\
\midrule
(1) & $H_0 = (\pi/2)\,c\,\xi_0^{10}/\bar\lambda_e$ & $h = 0{.}6682$ \\
(2) & $k_BT_0 = (16/9)\,\xi_0$ (Doc.~061) & $T_0 = 2{.}751$\,K \\
(3) & $\Omega_r \propto T_0^4/H_0^2$ & $9{.}643\times10^{-5}$ \\
(4) & $K_\text{frak} = 1-100\,\xi_0$ (A040) & $0{.}98667$ \\
(5) & $\int_0^\infty dz/E(z) = \pi/K_\text{frak}$ &
  $\boxed{\Omega_m^* = 0{.}3136}$ \\
\bottomrule
\end{tabular}
\end{center}
Planck: $\Omega_m = 0{.}315 \pm 0{.}007$ ($-0{.}19\,\sigma$).
Not a single step contains a parameter adjusted to $\Omega_m$.
Radiation ($\Omega_r \propto T_0^4/H_0^2$) is the link between
the particle and cosmic sectors.

\subsection{Overdetermination and open core}

\textbf{[K]} $T_0$ is determined via two independent routes:
particle sector (Route~1: $k_BT_0 = (16/9)\xi_0$) and cosmic sector
(Route~2: antipode condition). Both yield $\Omega_m$ agreeing to
$0{.}1\,\%$ --- the same overdetermination pattern as $\alpha$ in A130.

\textbf{[K]} Open (D(ii)): Coarse-grained entropy on the return half.
Poincaré discrepancy: $\sim 10^{10^{104}}$ dynamical times against
$\tau_c \approx 92$\,Gyr. Periodicity is a boundary condition, not a
consequence of the dynamics. Three cases (A050/Doc.~295):
Case~A ($d=0$, 1 cycle), Case~B ($\xi$ frozen, 75 cycles),
Case~C ($\xi$ running, log-spiral --- removes the entropy requirement,
but costs $\tau_c$).

\subsection{Falsification points}

\textbf{[K]} Low-$\ell$ cutoff: $\ell_\text{min} = 2\pi D_\text{LSS}/L^* = 3{.}08$
--- matches the quadrupole suppression unexplained since COBE (1992).
Primary falsification test: $m_\tau = 1776{.}97$\,MeV (Belle-II).

% ===================== 6 =====================
\section{Level 5 --- What separates the levels}

\textbf{[B]} The Landauer accounting (levels~2 and~3) does not touch
the $\xi$ scheme. It holds in any system with Hamiltonian or unitary
dynamics --- independent of the torus geometry. Explicitly booked in
A271: \enquote{Contact with the $\xi$ scheme of FFGFT: none.}

The arrow of time (level~1) is directly tied to $D_f = 3-\xi_0$ and
$\tilde{T} \cdot m = 1$ --- FFGFT-specific.
The cosmological thermodynamics (level~4) depends entirely on $\xi_0$,
$H_0$ and the torus geometry.

\begin{center}
\small
\begin{tabular}{@{}p{0.30\textwidth}p{0.20\textwidth}p{0.22\textwidth}p{0.18\textwidth}@{}}
\toprule
\textbf{Level} & \textbf{Document} & \textbf{$\xi$ dependence} & \textbf{Status} \\
\midrule
Arrow of time & Doc.~157 & direct & [B]/[S] \\
Landauer & A271 & none & [B] \\
Universal accounting & A280--A282 & none & [B] \\
Cosmological thermo. & Docs.~312/313 & direct & [K] \\
Information physics & Docs.~243--301 & indirect & [S]/[B] \\
\bottomrule
\end{tabular}
\end{center}

% ===================== 7 =====================
\section{Chain of argument: from geometry to heat}

\textbf{[B]} The five levels form a chain from FFGFT geometry to
classical thermodynamics:
(1) $\xi_0$ and $T^4$ fix the geometry.
(2) $\tilde{T} \cdot m = 1$ and the torus winding define a time
direction geometrically.
(3) In this directed time there is more phase-space volume in the
forward direction --- $\Delta S \ge 0$ as a geometric consequence.
(4) Non-bijective phase-space maps (Landauer) cost heat ---
independently of $\xi_0$, in the same phase-space picture.
(5) The thermal history ($\Omega_m^*$, $\Lambda^*$) follows from
$\xi_0$ directly.

Liouville's theorem and the second law are \emph{consequences}
in this chain, not postulates.

% ===================== 8 =====================
\section{Information thermodynamics in the IPI context (Docs.~243--301)}

\textbf{[S]} A further group of documents connects FFGFT thermodynamics
with the information-physics community (IPI context):
Doc.~251 (comparison with Vopson's information-mass hypothesis),
Doc.~247 (Landauer's principle in the context of information states),
Doc.~290 (information question: magnitude/phase),
Doc.~301 (information reversal test).
These documents are plausibility sketches or partial algebraic
derivations; they open connections to current debates without
contributing new core results to FFGFT thermodynamics.

% ===================== 9 =====================
\section{Phases, radiation and temperature in the time cycle (Doc.~313 in detail)}

\subsection{A gradient on a circle}

\textbf{[B]} The observable history shows a clear thermal sequence:
radiation domination, transition, matter domination. This gives the
impression of a directed process with a beginning. The FFGFT answer
(Doc.~313): the gradient is not a development \emph{from} a beginning
but a map \emph{along} the time cycle --- like climate zones along a
meridian circle.

\subsection{Three phases on the half-cycle}

\textbf{[K]} The observable history fills one half-cycle
($\tau_c/2 \approx 46$\,Gyr). Cycle positions:

\begin{center}
\small
\begin{tabular}{lccc}
\toprule
& & \multicolumn{2}{c}{$\theta/\pi$} \\
\cmidrule(l){3-4}
Event & $z$ & smooth & fractal \\
\midrule
Matter-radiation equality & $\sim\!3400$ & $0{.}93$ & $0{.}92$ \\
Scatter wall (recombination) & $1090$ & $0{.}992$ & $\mathbf{0{.}979}$ \\
Particle horizon & $z_\text{max}$ & $1{.}012$ & $\mathbf{0{.}9986}$ \\
\bottomrule
\end{tabular}
\end{center}

\subsubsection{Phase~1 --- Radiation domination (near the antipode)}

In the mass-flow reading: masses $m(\theta)$ are minimal, intrinsic
times $\tilde{T} = 1/m$ maximal --- clocks run slowest. The endpoint
slope of the mass profile at the antipode:
\[
  \left|\frac{dm}{d\theta}\right|_\text{antipode}
  = \sqrt{\Omega_r} \approx 0{.}0096
\]
No free parameter --- the observed radiation density appears here as
the \emph{smoothness condition} for periodic closure.

\subsubsection{Phase~2 --- Matter-radiation transition ($z \approx 3000$ to $1090$)}

With the fractal path correction $K_\text{frak} = 1-100\xi_0$ (A040,
independently derived) the particle horizon sits at $0{.}14\,\%$ from
the antipode --- improvement by factor~9 over the smooth calculation
($1{.}2\,\%$), no adjusted parameter.

The path correction applies to path/structure, cancels for path/path.
Invariance test: CMB peak positions unchanged (test passed without being
designed for it).

\subsubsection{Phase~3 --- Matter domination (today)}

Energy density $\propto (1+z)^3$. In the mass-flow reading masses
grow with decreasing $z$; clocks run faster; structure formation
unfolds.

\subsection{Temperature as a structural quantity}

\textbf{[K]} In the Standard Model: $T(z) = T_0(1+z)$ --- adiabatic
back-extrapolation. In FFGFT: $T_0$ comes from the atomic scale
(Doc.~061):
\[
  k_B T_0 = \tfrac{16}{9}\,\xi_0 \approx 2{.}370\times10^{-4}\,\text{eV}
  \quad(\text{measured: }2{.}349\times10^{-4}\,\text{eV, }{+0{.}92\,\%})
\]
The CMB temperature is a \emph{structural quantity} of the $\xi$ scheme,
not an endpoint of a cooling history.

\textbf{[K]} Range limit (Doc.~313): equilibria at other locations on
the time cycle are \emph{not} derivable from the current one by
extrapolation. Unaffected: the $\Omega_m^*$ chain (runs via $\xi_0$,
not via adiabatic back-calculation).

\subsection{Overdetermination of $T_0$}

\textbf{[K]} Two independent routes:
Route~1 (particle sector): $k_BT_0 = (16/9)\xi_0$, no cosmological input.
Route~2 (cosmic sector): antipode condition, no particle-physics input.
Both yield $\Omega_m$ agreeing to $0{.}1\,\%$
($0{.}3136$ vs.\ $0{.}3139$) --- the same pattern as $\alpha$ in A130.
Lever: $dT_0/T_0 \big/ d\Omega_m/\Omega_m \approx -12$.

\subsection{Lower bound and Hawking}

\textbf{[K]} Energy quantum of the time cycle (Doc.~312):
$E_\text{min} = \hbar H_0 = 2{.}28\times10^{-52}$\,J.
Finite $z_\text{max} = m_e/m_\text{min} - 1 = 3{.}59\times10^{38}$.
Singularity excluded: masses cannot approach zero.

\textbf{[K]} One KMS rule links Gibbons--Hawking, Unruh and Hawking:
$T = \hbar/(k_B\tau)$ with $\tau$ the local time period.
The compact time direction delivers $T_\text{GH}$ without a horizon.
The family ladder meets the cosmos at $r_s = R_H/2$ (exactly $0{.}500$).
Hawking evaporation for $M_\odot$: $56$ orders of magnitude too slow
for $\tau_c$ --- sharpens the entropy obstacle D(ii).

\subsection{Low-$\ell$ cutoff}

\textbf{[K]} On a torus of side length $L^*$ no modes with wavelength
$> L^*$ exist:
\[
  \ell_\text{min} = \frac{2\pi D_\text{LSS}}{L^*} = 3{.}08
\]
No adjusted parameter. Matches the quadrupole suppression
($\ell = 2$, factor $\sim 1/6$), unexplained since COBE (1992)
in $\Lambda$CDM. Order-of-magnitude statement [K$|$P20]; exact mode
summation on the torus still to be done.

\subsection{Reliable and unreliable statements}

\textbf{[K]} Three classes (Doc.~313):

\begin{center}
\small
\begin{tabular}{p{3.2cm}p{4cm}p{3cm}}
\toprule
Class & Examples & Reliability \\
\midrule
Counts & $|\text{Aut}(D4)| = 1152$, kissing number~24 & exact \\
Same-level ratios & $\alpha$, $m_\mu/m_e$, Koide~$Q$ &
  to $10^{-10}$ \\
With scale bridge & $G$, $H_0$, absolute masses & anchor needed \\
Number/duration quantities & $\eta = n_B/n_\gamma$, window durations &
  not transferable \\
\bottomrule
\end{tabular}
\end{center}

BBN observables ($Y_p$, D/H, ${}^7$Li): not a pure same-level ratio ---
the same applies to $\Lambda$CDM.

% ===================== 10 =====================
\section{Mass formation, fractal refinement and the illusion of expansion}

\subsection{What is absent at the antipode: metric sharpness}

\textbf{[K]} All physical lengths tied to mass scale as $1/m$:
Compton wavelength $\bar\lambda = \hbar/(mc)$,
Bohr radius $a_0 \propto 1/m_e$, intrinsic time $\tilde{T} = 1/m$.
At the antipode ($m \to m_\text{min}$) all run toward their maximum.
Without localisable massive objects \emph{metric sharpness} is absent.

\textbf{[B]} $m = 0$ is structurally excluded: from
$\tilde{T} \cdot m = 1$ with $T \ge \tau_0$ it follows that there is no
empty state, only a thinnest one (A060). Mass and temporal curvature are two
readings of the same structure --- neither produces the other.

\textbf{[K]} The quantification follows from the double structure of length bounds.
FFGFT has two cutoffs that follow from $\xi$ alone:
\begin{itemize}[leftmargin=1.4em,itemsep=0.1em]
  \item Lower length bound: $L_0 = \xi\,\ell_P$ (A060/A015)
        $\Rightarrow$ upper energy bound $E_\text{max} = E_P/\xi$
        (UV cutoff).
  \item Upper length bound: $L^* = 4\bar\lambda_e/\xi_0^{10}$ (Doc.~312)
        $\Rightarrow$ lower energy bound $E_\text{min} = \hbar H_0$
        (IR cutoff).
\end{itemize}
From $\tilde{T} \cdot m = 1$ and the IR bound:
\[
  m_\text{min} = \frac{\hbar H_0}{c^2} = 2{.}54\times10^{-69}\,\text{kg}
\]
This is no additional stipulation --- it is the cosmological side of the
same reciprocity $\tilde{T} \cdot m = 1$ that A060 describes on the particle
scale. A060 calls this range \enquote{unquantified} because the cosmological
scale lies outside its scope (A010); Docs.~312/313 quantify it. The metric
limiting case is thus fully determined: $m_\text{min}$ is a derived value,
not an open point.

\subsection{Photons: no freezing}

\textbf{[B]} Photons carry no rest mass --- $\tilde{T} \cdot m = 1$
does not apply to them directly. A photon at the antipode does not
freeze; it propagates at $c$. CMB photons lose energy geometrically
on the way to us (A265):
\[
  \frac{dE}{dx} = -\xi\,E
  \;\Rightarrow\;
  1 + z = e^{\xi x}
\]
Achromatic: all wavelengths stretched by the same factor.
From our perspective everything proceeds uniformly. The view back
is a view to a \emph{location} on the time cycle, not into the past.

\subsection{Expansion as fractal inward refinement}

\textbf{[B]} In the Standard Model: space expands outward, scale
factor $a(t)$ grows. In FFGFT there is no expanding space. As
$m(\theta)$ increases, all characteristic lengths ($\propto 1/m$)
become \emph{smaller}. Matter becomes more structured, more finely
resolved. This is \textbf{inward refinement}, not outward growth.

\textbf{[B]} $K_\text{frak} = 1-100\xi_0$ (A040): dimension deficit
$D_f = 3-\xi_0$, locally $\sim 0{.}003\,\%$, cumulatively over 17
orders of magnitude $\sim 100$.

\textbf{[B]} Recursion $\xi_{n+1} = \xi_n(1-100\xi_n)$ (A050), fixed
point: zero. Rolled up: winding closes after 75 cycles ($\tau_c$-cycle).
Unrolled: logarithmic spiral $\rho = 75\,e^{\beta/2\pi}$ --- what the
Standard Model calls expansion.

\textbf{[K]} Both descriptions are observationally equivalent (A265).
Rulers and clocks change in the same ratio; no local observer detects
the scaling. What we perceive as expansion is the contrast between our
current masses and the redshifted light from the antipode.

\textbf{[S]} Image: Mandelbrot zoom. Zooming inward reveals ever finer
structures, no outward growth. The antipode is the unzoomed view
(minimal differentiation, pure radiation). Our current position is
the deep-zoomed view (sharp atomic and stellar structures). The Standard
Model inverts the picture --- a choice of coordinates, not a physical
statement.

% ===================== 11 =====================
\section{Open questions}

\textbf{[H]} Three questions remain open:
(1) Minimal coherence cost for quantum coherence independent of error
correction quality (A271, §\,7.4).
(2) Continuous entropy flux as a function of clock rate, bit width and
logical irreversibility (A271, §\,8).
(3) Metric structure in the limit $m \to m_\text{min}$ --- whether a
regime of pure field geometry without localisable particle modes exists.

% ===================== 12 =====================
\section{Layer-status summary}

\begin{center}
\small
\begin{longtable}{@{}p{0.70\textwidth}l@{}}
\toprule
\textbf{Statement} & \textbf{Status} \\
\midrule
\endfirsthead
\toprule
\textbf{Statement} & \textbf{Status} \\
\midrule
\endhead
$\tilde{T} > 0$ forces geometric arrow of time & \textbf{[B]} \\
Torus winding: topologically protected direction & \textbf{[B]} \\
$D_f = 3-\xi_0$ breaks time-reversal symmetry & \textbf{[B]} \\
$\Delta S \ge 0$ as geometric consequence & \textbf{[B]} \\
Landauer: $Q \ge k_BT\ln 2$ from Liouville & \textbf{[B]} \\
Content neutrality of thermodynamics & \textbf{[B]} \\
Bennett's reversible logic & \textbf{[B]} \\
$\Omega_m^* = 0{.}3136$ from $\xi_0$ chain & \textbf{[K]} \\
$\Lambda^* = 1/R_H^2$ geometric carrier & \textbf{[K]} \\
Overdetermination $T_0$: two sectors to $0{.}1\,\%$ & \textbf{[K]} \\
Singularity excluded by $\hbar H_0$ & \textbf{[K]} \\
$\ell_\text{min} = 3{.}08$ (low-$\ell$ cutoff) & \textbf{[K]$|$P20} \\
Slowest clocks at the antipode & \textbf{[K]} \\
Compton wavelength, Bohr radius $\propto 1/m$ & \textbf{[B]} \\
Photons achromatically redshifted: $1+z = e^{\xi x}$ & \textbf{[B]} \\
Expansion = fractal inward refinement & \textbf{[B]} \\
Recursion: fixed point zero, spiral unrolled & \textbf{[B]} \\
Observational equivalence with $\Lambda$CDM & \textbf{[K]} \\
Landauer does not touch $\xi$ scheme & \textbf{[B]} \\
Coarse-grained closure D(ii): open & open \\
$m=0$ structurally excluded: no empty state (A060) & \textbf{[B]} \\
Double cutoff: $L_0 \Rightarrow E_\text{max}$, $L^* \Rightarrow E_\text{min}$ & \textbf{[K]} \\
$m_\text{min} = \hbar H_0/c^2$ from IR cutoff (Docs.~312/313) & \textbf{[K]} \\
Local field modes resolvable at thinnest point? & \textbf{[H]} \\
Past Hypothesis as geometric necessity & \textbf{[S]} \\
Mandelbrot zoom as image & \textbf{[S]} \\
\bottomrule
\end{longtable}
\end{center}

\textbf{Balance:} 14$\times$ \textbf{[B]} $\cdot$
8$\times$ \textbf{[K]} $\cdot$
3$\times$ \textbf{[S]} $\cdot$
2$\times$ \textbf{[H]} $\cdot$
1$\times$ open.

% ===================== 13 =====================
\section{Closing remark: two opposing models}

\textbf{[S]} An honest concession belongs at the end.

For us --- as observers with our current masses and clocks ---
our perspective naturally prevails. We think in terms of light travel
distances. We measure redshifts. We experience time as running forward.
This is real in the sense that it is what we measure, and our
instruments are calibrated to it.

The Standard Model ($\Lambda$CDM) takes exactly this perspective as its
starting point: expanding space, the Big Bang as a beginning,
redshift as spatial stretching, the CMB as a snapshot of an early epoch.
It is a model in excellent agreement with observations.

FFGFT proposes a different choice of coordinates: no expanding space,
no Big Bang as an event, redshift as geometric energy loss, the CMB
as a view to a location on the time cycle. This description also agrees
with the standard tests (A265, cosmological degeneracy).

\textbf{[B]} What distinguishes the two models:

\begin{itemize}[leftmargin=1.4em,itemsep=0.2em]
  \item $\Lambda$CDM generalises from the observer's perspective
        \emph{outward}: expanding space.
  \item FFGFT generalises from the geometry \emph{inward}:
        compact torus, fractal refinement.
\end{itemize}

Both are models. Both step beyond direct observation and add an
interpretive layer.

\textbf{[S]} What remains: two structurally opposite ways to describe
the same thing. In one everything expands outward, in the other
everything refines inward. In one time has a beginning, in the other
it is a closed cycle. In one the CMB is a window into the past, in
the other a window onto a location.

This is not a flaw --- it is an open question, stated precisely
enough to be decidable. The falsification tests ($m_\tau$,
$\ell_\text{min}$, $\Omega_m^*$ chain) will show whether the
geometric inward perspective of FFGFT carries more than the expanding
outward perspective --- or less.

% ===================== SOURCES =====================
\section{Source references}

\begin{description}[leftmargin=3.2em,style=sameline,itemsep=0.3em]
\small
\item[{Doc.~061}] Pascher, J., \emph{Temperature units and CMB in T0
theory}. Legacy corpus, 2025.
\item[{Doc.~157}] Pascher, J., \emph{The arrow of time in T0 theory}.
Legacy corpus, 2025.
\item[{Doc.~312}] Pascher, J., \emph{The closure scale}. August 2026.
\item[{Doc.~313}] Pascher, J., \emph{No beginning: time cycle and
$\Omega_m^*$}. August 2026.
\item[{A040}] Pascher, J., \emph{Fractal correction}. A-series, 2026.
\item[{A050}] Pascher, J., \emph{Recursion}. A-series, 2026.
\item[{A265}] Pascher, J., \emph{Redshift as a static geometry effect}.
A-series, 2026.
\item[{A271}] Pascher, J., \emph{Landauer and phase-space accounting}.
A-series, July 2026.
\item[{A280--A282}] Pascher, J., \emph{Phase-space accounting, carrier
and information, the computing sphere}. A-series, July 2026.
\end{description}

\section{Use of AI tools}

This document was produced with the assistance of AI tools.
The questions, stipulations and all substantive decisions originate
with the author; the tools formulate and organise.
Responsibility for content and errors rests entirely with the author.
Full statement in A010.
